\scanpage{598}

\mspage{7}

Prouvons enfin que $\unsure{\mathfrak m|m}=\mathfrak m_\Phi$. Comme $\mathfrak m_\Phi$ est un idéal maximal, il suffit de prouver

\[
\mathfrak m_\Phi\subset\mathfrak m\ .
\]

Or si $A\in\unsure{\mathcal P|\mathfrak P}(E)$ est telle que $A\notin\Phi$, \deleted{\unsure{ou $f$}} \unsure{i.e.|c.à.d.} $\unsure{f_A|\varphi_A}\in\mathfrak m_\Phi$ \deleted{\unsure{$=$}} (cor. 1), on a aussi $f_A\in\mathfrak m$. Donc il suffit de prouver le

\uline{Lemme}. Soit $\Phi\in\unsure{\mathscr F|\mathcal F}_0$. Alors $\mathfrak m_\Phi$ est \deleted{\unsure{le}} l'idéal engendré par les $f_A$ pour $A\in\mathcal P(E)\setminus\Phi$\inserted{, et \deleted{\unsure{même}} plus précisément, ces $f_A$ sont dans $\mathfrak m_\Phi$, et tt $f\in\mathfrak m_\Phi$ est de la forme $ff_A$, avec $A$ \unsure{au|et} \unsure{comme} \unsure{ci-dessus}}.

\deleted{Soit} En effet, $f\in\mathfrak m_\Phi$ signifie qu'il existe $B\in\Phi$ tel que $f\mid B=0$. Mais cela signifie aussi que $f=ff_A$, où $A=E\setminus A$, d'où $A\notin\Phi$. \deleted{Cela \unsure{montre} que $f$}

\inserted{Cela prouve la proposition 2, et en $=$ temps :}

\uline{Corollaire 2}. Tout anneau quotient intègre de $k_E=\prod_{\unsure{\xi|\zeta}\in E}k_\xi$ est un corps.

\deleted{\uline{\unsure{Corollaire 2}}} Notons \unsure{maintnt|maintenant} que si $\Phi\in\mathscr F$ est un filtre quelconque \deleted{de} $E$, alors $k_\Phi$ est un anneau quotient de $k_E$, donc

\diagram{\begin{tikzcd}[ampersand replacement=\&,row sep=0.1em,column sep=small]
\text{\unsure{(a)|(A)}}\quad\&\mathrm{Max}(k_\Phi)\arrow[r,hook]\&\mathrm{Max}(k_E)\ .\\
\&\&\begin{tikzpicture}[baseline=-0.5ex]\node[rotate=90,inner sep=0pt]{$\simeq$};\end{tikzpicture}\\
\&\&\mathscr F_0
\end{tikzcd}}

% slanted
\margin{\unsure{Notons} pour tt système d'anneaux $(k_\xi)_{\xi\in E}$ \uline{\unsure{non nuls|tous unit.}}, et $\Phi,\Psi\in\mathscr F$ \deleted{\unsure{sif}} ($\Psi$ plus fin que $\Phi$ \unsure{(i.e. $\Phi\subset\Psi$)}) sss $\mathfrak m_\Phi\subset\mathfrak m_\Psi$ (\unsure{suffira|suffit|supposer} le \unsure{premier} \unsure{filtre|ultrafiltre}, \unsure{soit} \unsure{surj.|sans} \illegible{1}).}

Ceci \unsure{posé|précisé}, on trouve

\uline{Corollaire 3} : \deleted{\unsure{Les}} L'inclusion précédente identifie $\mathrm{Max}(k_\Phi)$ à l'ens des ultrafiltres $\Psi$ sur $E$ qui sont \uline{plus fins} que $\Phi$.

En effet, si $\Psi$ est un tel ultrafiltre, on

% archive sheet number: 597
